% Figure captions (greyscale figures; numbers come from numbers_fda.tex)
\newcommand{\capFosr}{Average curves and where jump height matters. (A) Mean force in multiples of body weight (BW) over the \nTrials{} analysed trials (line) with $\pm$ 1 SD (shaded); the dotted line is body weight. Take-off runs from movement onset to take-off and landing from touch-down to the lowest centre-of-mass (COM) position, each rescaled to 0--100\%. (B) Jump-height coefficient from function-on-scalar regression (BW per SD of jump height, adjusted for age and body mass) with pointwise (dark) and simultaneous (light) 95\% bands from resampling athletes. Pale background shading marks the regions where both estimation approaches agree.}
\newcommand{\capSpm}{Top versus bottom jump-height tertile, one mean curve per athlete. (A) Group mean curves. (B) Difference (top minus bottom) with pointwise (dark) and simultaneous (light) 95\% bootstrap bands. Background shading marks clusters significant in the permutation test with family-wise error control.}
\newcommand{\capBand}{Jumps of almost equal height (\bandLo--\bandHi\,cm, \bandN{} trials) still differ in shape. Thin grey lines are individual trials, split by curve profile; the thick black line is the profile mean. Top: take-off. Bottom: landing.}
\newcommand{\capProfiles}{Curve profiles form a continuum. For each phase, the scatter plot shows the first two functional principal component (FPC) scores of every trial, with symbols for the two $k$-means profiles, and the histogram shows the distribution of the first score shaded by profile. There is no gap between the profiles and no second mode.}
